%==============================================================
\section{Thực nghiệm và thảo luận}
\label{sec:experiments}
%==============================================================

Chương này tổng hợp các thí nghiệm định lượng về routing và dựng thang mốc cho các phương pháp RL.
Trọng tâm không chỉ là các con số, mà là \textbf{phương pháp đo}: trong một môi trường mà nhiễu do
seed, do layout và do topology đều lớn, kết luận ``phương pháp của tôi tốt hơn'' rất dễ sai -- và sai
theo hướng mình muốn tin.

\subsection{Thiết lập thực nghiệm}

\begin{itemize}[noitemsep]
  \item \textbf{Phần mềm:} Python 3.12, Qiskit 2.5.2, \code{qiskit-aer}, \code{gymnasium}; môi
    trường quản lý bằng \code{uv} với file khoá phiên bản (\code{uv.lock}).
  \item \textbf{Mạch thử:} mạch ngẫu nhiên sinh bởi \code{bench.mach\_ngau\_nhien}, chỉ gồm cổng một
    qubit và \cx{}; ví dụ chính là mạch 9 qubit với khoảng 95 cổng \cx{}.
  \item \textbf{Topology:} \code{line(9)}, \code{ring(9)}, \code{grid(3,3)}, \code{full(9)}.
  \item \textbf{Router:} \code{BasicSwap}, \code{SabreSwap} với ba heuristic, \code{SabreLayout};
    các policy \code{RandomPolicy} và \code{GreedyPolicy} chạy trong \code{RoutingEnv}.
  \item \textbf{Metric:} số SWAP (đếm \emph{ngay sau} routing, vì sau translation SWAP đã bị phân rã
    thành CX) và depth; mọi kết quả đều qua \code{CheckMap}.
\end{itemize}

\paragraph{Bốn nguyên tắc của benchmark harness.} Module \code{bench.py} được thiết kế để chặn bốn
kiểu kết luận sai (Bảng~\ref{tab:harness}).

\begin{table}[h]
  \centering
  \caption{Nguyên tắc thiết kế của \code{bench.py}.}
  \label{tab:harness}
  \small
  \begin{tabularx}{\linewidth}{@{}XX@{}}
    \toprule
    \textbf{Nguyên tắc} & \textbf{Chặn lỗi gì} \\
    \midrule
    luôn chạy \code{CheckMap}, báo lỗi nếu không hợp lệ & agent ``thắng'' bằng cách bỏ qua cổng khó \\
    \textbf{tách} seed sinh mạch khỏi seed của router & lẫn ``mạch này khó hơn'' với ``router dao động'' \\
    cố định layout khi so sánh router & so sai bài toán (layout ảnh hưởng $\sim$16\%, lớn hơn khác biệt giữa router) \\
    so sánh \textbf{ghép cặp} theo (topology, mạch, seed) & độ khó giữa các mạch lấn át khác biệt giữa router \\
    \bottomrule
  \end{tabularx}
\end{table}

\subsection{Ảnh hưởng của topology}
\label{sec:exp-topology}

Hình~\ref{fig:topology} cho thấy topology ảnh hưởng tới số SWAP mạnh hơn mọi lựa chọn heuristic. Chỉ
thêm \emph{một} cạnh (từ \code{line(9)} sang \code{ring(9)}) đã giảm số SWAP khoảng 44\%;
\code{full(9)} không cần SWAP nào -- routing chỉ tồn tại vì phần cứng thiếu kết nối. Đáng chú ý,
\code{grid(3,3)} cần ít SWAP hơn \code{ring(9)} nhưng depth lại cao hơn (227 so với 183 trong lần
chạy kiểm chứng): \textbf{số SWAP và depth không cùng chiều}, và đó là lý do phải chọn reward có ý thức
(Bảng~\ref{tab:reward}).

\begin{figure}[h]
  \centering
  \begin{tikzpicture}
    \begin{axis}[
      ybar, width=0.62\linewidth, height=5.2cm, bar width=18pt,
      ylabel={Số SWAP trung bình}, ymin=0, ymax=120,
      symbolic x coords={line,ring,grid,full},
      xtick=data, xticklabels={line(9),ring(9),grid(3{,}3),full(9)}, xticklabel style={font=\small\ttfamily},
      nodes near coords, nodes near coords style={font=\footnotesize, anchor=south west, xshift=3pt},
      grid=major, grid style={black!10}, axis x line*=bottom, axis y line*=left,
      enlarge x limits=0.18,
    ]
      \addplot[fill=AccentBlue!60, draw=AccentBlue,
               error bars/.cd, y dir=both, y explicit]
        coordinates {
          (line,75.2) +- (0,32.8)
          (ring,42.8) +- (0,15.8)
          (grid,24.5) +- (0,11.4)
          (full,0) +- (0,0)
        };
    \end{axis}
  \end{tikzpicture}
  \caption{Số SWAP khi route mạch 9 qubit ($\sim$95 \cx{}) bằng \code{SabreSwap(decay)} trên layout
  tầm thường, trung bình trên 5 mạch $\times$ 5 seed $=$ 25 lần đo mỗi topology. Thanh lỗi là độ lệch
  chuẩn, lớn vì trộn cả độ khó khác nhau giữa các mạch.}
  \label{fig:topology}
\end{figure}

\subsection{Nhiễu do seed và vai trò của layout}

SABRE phá thế hoà giữa các SWAP cùng cost bằng số ngẫu nhiên, nên kết quả phụ thuộc \code{seed}.
Tách riêng phần nhiễu này bằng \code{bench.nhieu\_router} (cố định mạch, chỉ đổi seed router) cho
thấy \code{SabreSwap(decay)} dao động khoảng \textbf{9\%}; \code{BasicSwap} thì tất định. Một kết
quả quan trọng cho thiết kế thí nghiệm:
\begin{keybox}
Một cải thiện nhỏ hơn mức dao động do seed ($\sim$9\%) là nhiễu, trừ khi được đo trên đủ nhiều seed để
độ lệch chuẩn của trung bình nhỏ lại. Lỗi phương pháp phổ biến nhất là so agent (chạy nhiều lần, lấy
kết quả tốt nhất) với baseline chạy \emph{một} seed.
\end{keybox}

Trên một mạch 9 qubit với 158 cổng \cx{}, 15 seed của \code{SabreSwap} trên layout tầm thường cho
$134.8 \pm 5.7$ SWAP (khoảng 127--149), còn \code{SabreLayout} -- chọn layout và route cùng lúc -- cho
$112.5 \pm 3.3$ SWAP. Trên tập benchmark, \code{SabreLayout} tốt hơn \code{SabreSwap} + layout tầm
thường khoảng \textbf{16\%}. Layout và routing không tách rời: một agent chỉ làm routing trên layout
cố định phải được so với \code{SabreSwap} \emph{trên cùng layout đó}; so với \code{SabreLayout} là so
sai bài toán.

Ngay cả khi tổng hợp trên 48 cặp (topology, mạch, seed) cho thấy \code{decay} tốt hơn \code{basic}
khoảng 10\%, vẫn có những cặp mà \code{basic} thắng. Chọn báo cáo đúng một cặp như vậy -- sau khi đã
xem kết quả -- sẽ ``chứng minh'' được kết luận ngược lại. Đây chính là hình dạng của một bảng số liệu
RL chỉ chạy một seed trên một mạch.

\subsection{Thang mốc cho reinforcement learning}

Bảng~\ref{tab:baselines} là thang mốc mà các agent ở Chuyên đề nghiên cứu 2 sẽ thêm một dòng vào: 20
mạch 5 qubit trên \code{line(5)}, layout tầm thường (vì môi trường bắt đầu từ đó), so sánh ghép cặp
từng mạch với \code{SabreSwap(decay)}.

\begin{table}[h]
  \centering
  \caption{Thang mốc trên \code{line(5)}, 20 mạch 5 qubit. Cột cuối là mức kém hơn
  \code{SabreSwap(decay)}, tính ghép cặp theo từng mạch. Số trong ngoặc là kết quả chạy lại để
  kiểm chứng.}
  \label{tab:baselines}
  \small
  \begin{tabular}{@{}lcc@{}}
    \toprule
    \textbf{Phương pháp} & \textbf{SWAP trung bình} & \textbf{Kém \code{SabreSwap(decay)}} \\
    \midrule
    \code{SabreSwap(decay)} -- mốc trên & 8.25 \;(8.30) & 0\% \\
    \code{SabreSwap(basic)} & 8.45 \;(9.15) & 2.3\% \;(11.2\%) \\
    \code{GreedyPolicy} trong \code{RoutingEnv} -- mốc giữa & 10.45 \;(10.45) & \textbf{24.9\%} \;(24.4\%) \\
    \code{BasicSwap} & 13.90 \;(13.90) & 64.6\% \;(63.7\%) \\
    \code{RandomPolicy} trong \code{RoutingEnv} -- mốc dưới & 27.95 \;(27.95) & 266\% \;(264.5\%) \\
    \bottomrule
  \end{tabular}
\end{table}

\begin{figure}[h]
  \centering
  \begin{tikzpicture}
    \begin{axis}[
      xbar, width=0.7\linewidth, height=4.6cm, bar width=9pt,
      xlabel={Số SWAP trung bình (\code{line(5)}, 20 mạch)}, xmin=0, xmax=33,
      symbolic y coords={rnd,bsw,grd,sbb,sbd},
      ytick=data, yticklabels={RandomPolicy,BasicSwap,GreedyPolicy,SabreSwap(basic),SabreSwap(decay)}, yticklabel style={font=\small\ttfamily},
      nodes near coords, nodes near coords style={font=\footnotesize},
      grid=major, grid style={black!10}, axis x line*=bottom, axis y line*=left,
      enlarge y limits=0.15,
    ]
      \addplot[fill=OkGreen!50, draw=OkGreen] coordinates {
        (27.95,rnd) (13.90,bsw) (10.45,grd) (8.45,sbb) (8.25,sbd)
      };
    \end{axis}
  \end{tikzpicture}
  \caption{Thang mốc cho agent RL. Khoảng cách 24.9\% giữa \code{GreedyPolicy} và
  \code{SabreSwap(decay)} là mục tiêu đầu tiên cần thu hẹp.}
  \label{fig:baselines}
\end{figure}

Số chạy lại khớp với số gốc ở mọi dòng trừ \code{SabreSwap(basic)} -- dòng duy nhất mà kết quả phụ
thuộc chi tiết cài đặt ngẫu nhiên bên trong SABRE, nhất quán với mức nhiễu do seed ở trên. Mục tiêu
đầu tiên của một agent \textbf{không phải} thắng \code{SabreSwap} mà là \textbf{vượt}
\code{GreedyPolicy}, tức thu hẹp khoảng 25\%. Chỉ hơn \code{RandomPolicy} thì chưa chứng minh được gì.

\subsection{Năm ``bẫy im lặng'' và bài học phương pháp}

Trong quá trình thực nghiệm đã gặp năm lỗi không sinh ra exception nào, chỉ sinh ra số liệu sai.
Chúng đều đánh trúng người tự dựng pipeline routing, nên được ghi lại đầy đủ (Bảng~\ref{tab:traps}).

\begin{table}[h]
  \centering
  \caption{Năm bẫy im lặng gặp trong thực nghiệm và cách phát hiện.}
  \label{tab:traps}
  \small
  \begin{tabularx}{\linewidth}{@{}lXX@{}}
    \toprule
    & \textbf{Bẫy} & \textbf{Cách phát hiện / tránh} \\
    \midrule
    1 & \code{SabreLayout} tự route; nối với \code{ApplyLayout} thì áp layout hai lần $\to$ mạch không hợp lệ & pass nào ghi \code{final\_layout} thì dùng một mình; luôn chạy \code{CheckMap} \\
    2 & \code{Operator(compiled).equiv(Operator(goc))} trả \code{False} dù mạch đúng & dùng \code{Operator.from\_circuit(compiled)} \\
    3 & bù hoán vị bằng $\pi$ thay vì $\pi^{-1}$; chỉ lộ khi hoán vị không phải một phép đổi chỗ & test với hoán vị phức tạp; bù bằng nghịch đảo \\
    4 & \code{unitary\_synthesis\_method} chỉ tác dụng lên \code{UnitaryGate}; \code{epsilon} bị bỏ qua với \code{rz} thường & luôn đo sai số thực của mạch tổng hợp \\
    5 & plugin Ross--Selinger tên \code{"gridsynth"}; plugin \code{"sk"} gây panic ở tầng Rust, không bắt được bằng \code{try/except} & dùng pass \code{SolovayKitaev} thay cho plugin \\
    \bottomrule
  \end{tabularx}
\end{table}

\paragraph{Ba điều dễ hiểu sai về RL cho biên dịch.}
\begin{enumerate}[noitemsep]
  \item \textbf{Hợp lệ không đồng nghĩa với đúng.} \code{CheckMap} chỉ kiểm tra cổng hai qubit nằm
    trên cạnh thật; một agent vẫn có thể sinh mạch hợp lệ mà sai unitary. Phải kiểm tra cả hai, ở mọi
    episode.
  \item \textbf{Baseline yếu làm kết quả vô nghĩa.} \code{mini\_router} cài đúng công thức của SABRE
    mà vẫn tốn 10 SWAP so với 8 của \code{SabreSwap}. So agent với baseline tự cài thì agent ``thắng''
    dễ dàng, và con số đó không nói lên điều gì.
  \item \textbf{Ít SWAP không đồng nghĩa với mạch tốt hơn.} Các reward \code{swap}, \code{depth},
    \code{fidelity} không cùng chiều; \code{fidelity} gần mục tiêu cuối cùng nhất.
\end{enumerate}
